% ══════════════════════════════════════════════════════════════════════════════
% APPENDIX A1 — DICTIONARY AND BLOCKLIST  [v31 final]
% ══════════════════════════════════════════════════════════════════════════════

This appendix documents the rule-based components of the text-measurement
pipeline. It gives the final dictionary used to construct \netscore{} and the
Random Forest blocklist used to remove identity and boilerplate leakage. The
purpose is reproducibility: every dictionary score and every blocklisted RF
feature can be traced to the exact term lists below.

\section{Dictionary design}
\label{app:a1_dictionary_design}

The dictionary is organised around four semantic categories: policy-action
language, inflation language, employment and real-activity language, and a
residual category for risk, crisis, and commitment terms. The word ``residual''
is used only as a dictionary-category label; it is not an econometric residual
and should be read as a catch-all group for stance-relevant terms that do not
fit cleanly into policy, inflation, or employment/growth language. Each term is assigned
a direction, hawkish or dovish. The document-level score pools hawkish and
dovish matches and applies the Laplace-smoothed formula described in
Chapter~\ref{ch:methodology}:
\[
  \netscore_i = \frac{H_i-D_i}{H_i+D_i+10}.
\]
The same formula is used for category-level sub-scores. The smoothing constant
shrinks documents with few stance-bearing matches toward zero and prevents one
or two isolated terms from generating an extreme score.

All dictionary terms are matched after conservative stop-word removal and
Porter stemming. The term lists below are shown in their human-readable form,
not in their stemmed form. The scoring pipeline stems both the document text and
the dictionary entries before matching. The final version contains 288 raw phrase
entries, split almost evenly between hawkish and dovish terms. After stemming
and deduplication, these correspond to 282 unique bigram/trigram patterns.

\section{Stop-word treatment}
\label{app:a1_stopwords}

The custom stop-word list removes only non-directional function words. Words
such as \emph{above}, \emph{below}, \emph{further}, \emph{more}, \emph{not},
and \emph{no} are deliberately retained because they can determine the stance
of a monetary-policy phrase.

\begin{table}[H]
\centering
\caption{Custom stop-word list used in dictionary scoring}
\label{tab:a1_stopwords}
\small
\begin{tabular}{llllll}
\toprule
\multicolumn{6}{c}{Non-directional function words removed before dictionary matching} \\
\midrule
a & an & the & is & are & was \\
were & be & been & being & has & have \\
had & having & do & does & did & it \\
its & our & their & we & they & i \\
me & my & he & she & him & her \\
his & of & to & at & by & with \\
on & in & for & and & or & but \\
that & which &  &  &  &  \\
\bottomrule
\end{tabular}
\end{table}


\section{Final dictionary terms}
\label{app:a1_terms}

\begin{longtable}{p{0.47\textwidth}p{0.47\textwidth}}
\caption{Policy-action dictionary terms}\label{tab:a1_policy_terms}\\
\toprule
Hawkish terms & Dovish terms \\
\midrule
\endfirsthead
\toprule
Hawkish terms & Dovish terms \\
\midrule
\endhead
decided raise & lower target \\
raise interest & lower rate \\
raise interest rates & lower interest rates \\
raising interest rates & lowering interest rates \\
raised interest rates & decided lower \\
raise target & decided reduce \\
raising rates & reduce key \\
decided increase & cut rates \\
increase rates & point reduction \\
increasing rates & point decrease \\
point increase & rate cut \\
ongoing increases & point cut \\
rate increase & rate reduction \\
rate hike & appropriate lower \\
raise rates & appropriate reduce \\
ongoing rate increases & pace reductions \\
further rate increases & slower pace \\
pace increases & ease monetary policy \\
further increases & policy easing \\
additional increases & monetary easing \\
sufficiently restrictive & further easing \\
restrictive policy & policy accommodation \\
restrictive stance & monetary accommodation \\
restrictive territory & accommodative policy \\
restrictive monetary policy & accommodative stance \\
restrictive policy stance & highly accommodative \\
more restrictive & remains accommodative \\
policy tightening & ample degree \\
monetary tightening & asset purchases \\
further tightening & purchase programme \\
tighten monetary policy & purchase program \\
tightening monetary policy & net asset purchases \\
additional policy firming & net purchases \\
further policy firming & quantitative easing \\
cumulative tightening & forward guidance \\
quantitative tightening & reinvesting principal \\
balance sheet reduction & reinvestment policy \\
tightening cycle & targeted longer \\
remove accommodation & pandemic emergency \\
faster pace & favourable financing \\
restore price stability & preserve favourable \\
bring inflation down & ample liquidity \\
bring inflation back & smooth transmission \\
returning inflation & support economy \\
appropriate raise & support recovery \\
remain vigilant & support economic \\
vigilant about & extended period \\
maintain price stability & least until \\
strongly committed & until substantial further \\
firmly committed & lower bound \\
policy normalization & neutral rate \\
normalize policy & monetary stimulus \\
committed returning & policy stimulus \\
 & additional stimulus \\
 & further stimulus \\
 & generate economic \\
 & stabilises sustainably \\
\bottomrule
\end{longtable}

\begin{longtable}{p{0.47\textwidth}p{0.47\textwidth}}
\caption{Inflation dictionary terms}\label{tab:a1_inflation_terms}\\
\toprule
Hawkish terms & Dovish terms \\
\midrule
\endfirsthead
\toprule
Hawkish terms & Dovish terms \\
\midrule
\endhead
high inflation & low inflation \\
higher inflation & lower inflation \\
rising inflation & subdued inflation \\
inflation remains elevated & muted inflation \\
inflation high & falling inflation \\
inflation elevated & inflation declined \\
elevated inflation & inflation fell \\
inflation remains high & inflation dropped \\
still too high & inflation eased \\
inflation running & inflation fallen \\
far too high & continued decline \\
remain well above & moving down \\
significantly above & moving sustainably toward \\
above target & ongoing disinflation \\
considerably above & disinflation process \\
higher prices & disinflationary pressures \\
rising prices & prices fell \\
high prices & projected decline \\
prices rose & lower prices \\
inflation rose & transitory factors \\
price pressures & temporary factors \\
price pressures remain & below target \\
core inflation & below aim \\
underlying inflation & significantly below \\
services inflation & price stability objective \\
domestic price & closer target \\
producer prices & back target \\
upward trend & deflation risk \\
wage growth & deflationary pressure \\
wage pressures & deflationary risk \\
nominal wage & well anchored \\
labour costs & firmly anchored \\
unit labour cost &  \\
unit labor cost &  \\
second round effects &  \\
wage price spiral &  \\
cost pressures &  \\
pass through &  \\
tight labor &  \\
tight labour &  \\
upward pressure &  \\
inflationary pressures &  \\
inflationary risks &  \\
inflationary expectations &  \\
entrenched inflation &  \\
inflation overshoot &  \\
pricing power &  \\
\bottomrule
\end{longtable}

\begin{longtable}{p{0.47\textwidth}p{0.47\textwidth}}
\caption{Employment and real-activity dictionary terms}\label{tab:a1_employment_terms}\\
\toprule
Hawkish terms & Dovish terms \\
\midrule
\endfirsthead
\toprule
Hawkish terms & Dovish terms \\
\midrule
\endhead
strong growth & economic weakness \\
robust growth & weak growth \\
higher growth & low growth \\
high growth & lower growth \\
rapid growth & weaker growth \\
above potential & slower growth \\
solid growth & modest growth \\
solid pace & weak demand \\
strong pace & weaker demand \\
strong economic & lower demand \\
strong demand & growth slowed \\
economic expansion & drag growth \\
continued expand & economic slowdown \\
strong labor market & economic downturn \\
strong labour market & weak economic \\
tight labor market & weaker economic \\
tight labour market & sharp decline \\
very tight & negative growth \\
very strong labor & deep recession \\
low unemployment & high unemployment \\
job gains & rising unemployment \\
rising wages & higher unemployment \\
higher wages & labor market slack \\
strong wages & job losses \\
supply constraints & loss momentum \\
capacity constraints &  \\
labour shortage &  \\
labor shortage &  \\
\bottomrule
\end{longtable}

\begin{longtable}{p{0.47\textwidth}p{0.47\textwidth}}
\caption{Residual risk and commitment dictionary terms}\label{tab:a1_residual_terms}\\
\toprule
Hawkish terms & Dovish terms \\
\midrule
\endfirsthead
\toprule
Hawkish terms & Dovish terms \\
\midrule
\endhead
upside risk & downside risk \\
tail risks & risks growth \\
strong vigilance & path inflation \\
continued vigilance & negative territory \\
remain alert & economic headwinds \\
closely monitor & credit tightening \\
upward revision & tighter credit \\
revised upward & tighter financing \\
revised up & tightening financial \\
stronger than expected & credit standards \\
higher than expected & market turmoil \\
supply shock & systemic risk \\
supply disruption & financial vulnerabilities \\
energy shock & liquidity trap \\
second round & zero lower bound \\
 & effective lower bound \\
 & downward revision \\
 & revised downward \\
 & revised down \\
 & weaker than expected \\
 & lower than expected \\
 & geopolitical risks \\
 & geopolitical tensions \\
 & trade tensions \\
 & sovereign debt \\
 & debt crisis \\
 & fiscal consolidation \\
 & high uncertainty \\
 & heightened uncertainty \\
 & elevated uncertainty \\
 & financial fragmentation \\
\bottomrule
\end{longtable}


\section{Random Forest blocklist}
\label{app:a1_blocklist}

The Random Forest uses a separate blocklist before TF--IDF construction. The
blocklist is not a tone dictionary. Its purpose is to remove bigrams that could
predict the rate-decision class through speaker identity, institutional
boilerplate, location names, or period-specific artefacts rather than through
genuine stance language. This is important because central bank leaders often
serve during particular policy regimes: a chair or president name can become a
spurious proxy for a tightening or easing cycle.

\begin{longtable}{p{0.47\textwidth}p{0.47\textwidth}}
\caption{Random Forest blocklist bigrams}\label{tab:a1_blocklist}\\
\toprule
Boilerplate and institutional bigrams & Names, places, and era markers \\
\midrule
\endfirsthead
\toprule
Boilerplate and institutional bigrams & Names, places, and era markers \\
\midrule
\endhead
cleveland richmond & alan greenspan \\
richmond atlanta & greenspan chairman \\
atlanta boston & chairman greenspan \\
chicago st & ben bernank \\
chicago minneapoli & bernank chairman \\
loui minneapoli & chairman bernank \\
st loui & janet yellen \\
kansa citi & yellen chairman \\
san francisco & chairman yellen \\
new york & jerom powel \\
reserv bank & powel chairman \\
feder reserv & chairman powel \\
york philadelphia & chairman susan \\
philadelphia cleveland & susan bie \\
cleveland atlanta & bie roger \\
atlanta chicago & roger fergus \\
minneapoli kansa & frederic mishkin \\
dalla san & donald kohn \\
discount rate & randal kroszner \\
request submit & kroszner freder \\
board approv & kevin warsh \\
relat action & wim duisenberg \\
committe continu & jean claud \\
submit board & claud trichet \\
approv request & mario draghi \\
governor approv & christin lagard \\
request discount & lagard christin \\
charg depositori & luca papademo \\
depositori institut & vitor constanci \\
press confer & constanci vitor \\
confer start & de guindo \\
relat content & philip lane \\
consider under & isabel schnabel \\
polici measur & piero cipollon \\
polici decis & washington dc \\
press releas & constitut avenu \\
monetari polici & 20th street \\
polici transmiss & frankfurt main \\
communic reproduct & sonnemannstrass 20 \\
reproduct permit & european central \\
permit provid & central bank \\
council steer & centuri low \\
pre commit & new centuri \\
minimum bid &  \\
bid rate &  \\
futur roll &  \\
avoid interfer &  \\
time return &  \\
american famili &  \\
social partner &  \\
mandat goal &  \\
\bottomrule
\end{longtable}


\section{Interpretation}
\label{app:a1_interpretation}

The dictionary is intentionally transparent and static. This makes it easy to
interpret but also creates two limitations. First, the dictionary cannot fully
resolve context, negation, or changes in meaning over time. Second, the same
vocabulary is applied to both the FOMC and the ECB even though their mandates
and communication styles differ. These limitations are why the thesis compares
\netscore{} against the Random Forest and RoBERTa/BERT scores throughout
Chapter~\ref{ch:results}.


% ══════════════════════════════════════════════════════════════════════════════
% APPENDIX — Top 10 dictionary terms per sub-category (grayscale)
% Requires: booktabs, threeparttable, tabularx, xcolor (for row shading)
% ══════════════════════════════════════════════════════════════════════════════

\begin{table}[htbp]
\centering
\caption[Top dictionary terms by sub-category]{Top 10 dictionary terms per sub-category ranked by total corpus hits (deduplicated by stemmed form). FOMC counts pool statements and press conferences (334 documents); ECB counts pool press releases and press conferences (550 documents).}
\label{tab:app_top_dict_terms}
\begin{threeparttable}
\scriptsize
\setlength{\tabcolsep}{2.0pt}
\renewcommand{\arraystretch}{1.02}

\vspace{0.25cm}
\noindent
\begin{tabularx}{\textwidth}{@{}>{\itshape}Xrr|>{\itshape}Xrr@{}}
\toprule
\multicolumn{3}{@{}l}{\textbf{Policy --- Hawkish}} & \multicolumn{3}{l@{}}{\textbf{Policy --- Dovish}} \\
\cmidrule(r){1-3}\cmidrule(l){4-6}
\normalfont Term & \normalfont FOMC & \normalfont ECB & \normalfont Term & \normalfont FOMC & \normalfont ECB \\
\midrule
rate increase & 160 & 204 & asset purchases & 400 & 655 \\
rate hike & 146 & 169 & forward guidance & 119 & 396 \\
raise rates & 206 & 70 & rate cut & 172 & 305 \\
returning inflation & 106 & 130 & purchase programme & 0 & 474 \\
raise interest & 100 & 107 & extended period & 59 & 209 \\
raise interest rates & 99 & 107 & net asset purchases & 11 & 247 \\
maintain price stability & 17 & 163 & favourable financing & 0 & 227 \\
increase rates & 23 & 150 & cut rates & 118 & 89 \\
sufficiently restrictive & 57 & 71 & monetary accommodation & 9 & 193 \\
strongly committed & 101 & 6 & support economic & 76 & 123 \\
\bottomrule
\end{tabularx}

\vspace{0.25cm}
\noindent
\begin{tabularx}{\textwidth}{@{}>{\itshape}Xrr|>{\itshape}Xrr@{}}
\toprule
\multicolumn{3}{@{}l}{\textbf{Inflation --- Hawkish}} & \multicolumn{3}{l@{}}{\textbf{Inflation --- Dovish}} \\
\cmidrule(r){1-3}\cmidrule(l){4-6}
\normalfont Term & \normalfont FOMC & \normalfont ECB & \normalfont Term & \normalfont FOMC & \normalfont ECB \\
\midrule
underlying inflation & 44 & 390 & low inflation & 75 & 162 \\
inflationary pressures & 46 & 309 & firmly anchored & 1 & 206 \\
second round effects & 3 & 294 & well anchored & 154 & 43 \\
core inflation & 140 & 126 & moving down & 118 & 7 \\
price pressures & 21 & 242 & lower inflation & 30 & 37 \\
wage growth & 113 & 133 & inflation declined & 17 & 41 \\
high inflation & 133 & 64 & continued decline & 14 & 43 \\
upward pressure & 68 & 116 & temporary factors & 8 & 47 \\
pass through & 34 & 145 & inflation eased & 35 & 11 \\
labour costs & 0 & 173 & below target & 22 & 21 \\
\bottomrule
\end{tabularx}

\vspace{0.25cm}
\noindent
\begin{tabularx}{\textwidth}{@{}>{\itshape}Xrr|>{\itshape}Xrr@{}}
\toprule
\multicolumn{3}{@{}l}{\textbf{Employment --- Hawkish}} & \multicolumn{3}{l@{}}{\textbf{Employment --- Dovish}} \\
\cmidrule(r){1-3}\cmidrule(l){4-6}
\normalfont Term & \normalfont FOMC & \normalfont ECB & \normalfont Term & \normalfont FOMC & \normalfont ECB \\
\midrule
strong labor market & 175 & 0 & high unemployment & 25 & 35 \\
job gains & 174 & 0 & growth slowed & 37 & 16 \\
strong growth & 34 & 77 & slower growth & 24 & 22 \\
continued expand & 67 & 42 & weak economic & 2 & 33 \\
economic expansion & 32 & 71 & sharp decline & 9 & 20 \\
low unemployment & 54 & 6 & weak demand & 5 & 24 \\
solid pace & 59 & 0 & job losses & 28 & 0 \\
supply constraints & 24 & 27 & lower growth & 5 & 23 \\
rising wages & 11 & 39 & economic weakness & 16 & 11 \\
tight labor market & 50 & 0 & weaker growth & 6 & 21 \\
\bottomrule
\end{tabularx}

\vspace{0.25cm}
\noindent
\begin{tabularx}{\textwidth}{@{}>{\itshape}Xrr|>{\itshape}Xrr@{}}
\toprule
\multicolumn{3}{@{}l}{\textbf{Residual --- Hawkish}} & \multicolumn{3}{l@{}}{\textbf{Residual --- Dovish}} \\
\cmidrule(r){1-3}\cmidrule(l){4-6}
\normalfont Term & \normalfont FOMC & \normalfont ECB & \normalfont Term & \normalfont FOMC & \normalfont ECB \\
\midrule
upside risk & 31 & 383 & downside risk & 147 & 353 \\
second round & 8 & 297 & fiscal consolidation & 5 & 294 \\
closely monitor & 54 & 87 & path inflation & 11 & 137 \\
strong vigilance & 0 & 90 & risks growth & 12 & 89 \\
stronger than expected & 15 & 60 & credit standards & 6 & 87 \\
higher than expected & 12 & 52 & high uncertainty & 9 & 78 \\
revised up & 17 & 47 & geopolitical risks & 11 & 67 \\
revised upward & 0 & 50 & downward revision & 13 & 51 \\
tail risks & 8 & 35 & weaker than expected & 6 & 58 \\
upward revision & 5 & 31 & sovereign debt & 9 & 49 \\
\bottomrule
\end{tabularx}

\begin{tablenotes}[flushleft]
\footnotesize
\item \emph{Notes}: Counts are raw term matches after Porter stemming and
stop-word removal. Where multiple raw phrases map to the same stemmed form
(e.g.\ ``raise rates'' and ``raising rates''), only the shortest form is shown
and counts are reported once. A single document can contribute multiple hits
for the same term. The full dictionary (288 terms) is reported in
Appendix~\ref{app:a1_terms}.
\end{tablenotes}
\end{threeparttable}
\end{table}